\documentclass[12pt,a4paper]{article}
\usepackage[utf8]{inputenc}
\usepackage[T1]{fontenc}
\usepackage[french]{babel}
\usepackage{geometry}
\usepackage{setspace}
\usepackage{lmodern}
\usepackage{microtype}
\usepackage{amsmath,amssymb}
\usepackage{booktabs}
\usepackage{array}
\usepackage{enumitem}
\usepackage{hyperref}
\usepackage{tikz}
\usetikzlibrary{arrows.meta,positioning,automata,shapes.geometric}

\geometry{left=2.5cm,right=2.5cm,top=2.5cm,bottom=2.5cm}
\setstretch{1.15}
\setlist[itemize]{leftmargin=1.2cm}
\setlist[enumerate]{leftmargin=1.2cm}
\hypersetup{colorlinks=true,linkcolor=black,urlcolor=blue,citecolor=black}

\begin{document}

% =========================================
% PAGE DE GARDE
% =========================================
\begin{titlepage}
    \centering

    {\Large \textbf{Universit\'e M'Hamed Bougara de Boumerd\`es (UMBB)}}\\[0.3cm]
    {\large D\'epartement d'Informatique}\\[0.8cm]
    {\large \textbf{Syst\`emes Multi-Agents}}\\[0.2cm]
    {\large \textbf{M1 G\'enie Logiciel}}\\[1.8cm]

    {\LARGE \textbf{Rapport Technique}}\\[0.3cm]
    {\Large Projet SRUU : Syst\`eme de R\'eponse d'Urgence Urbaine}\\[1.6cm]

    \begin{tabular}{>{\bfseries}p{4.5cm}p{8.8cm}}
        R\'ealis\'e par & Djemil Hana \\
        Groupe & 2 \\
        Encadrante & Refas Ikram \\
        Ann\'ee universitaire & 2025--2026 \\
    \end{tabular}

    \vfill
    {\large \today}
\end{titlepage}

\tableofcontents
\newpage

% =========================================
\section{Introduction}
Ce rapport pr\'esente les choix d'architecture et de conception du syst\`eme \textbf{SRUU} (Syst\`eme de R\'eponse d'Urgence Urbaine), impl\'ement\'e avec la plateforme JADE. L'objectif est de coordonner plusieurs types d'agents pour d\'etecter, prioriser et traiter des incidents urbains h\'et\'erog\`enes (incendie, urgence m\'edicale, effondrement structurel, biohazard, fuite cryog\'enique).

L'approche retenue combine :
\begin{itemize}
    \item une architecture multi-agents distribu\'ee ;
    \item un protocole d'allocation de ressources bas\'e sur \textbf{FIPA Contract Net} ;
    \item une ontologie commune pour garantir la coh\'erence s\'emantique ;
    \item une d\'ecision du r\'epartiteur par fonction d'utilit\'e multi-crit\`ere ;
    \item des comportements agents structur\'es en machines \`a \`etats finis (FSM) ;
    \item des m\'ecanismes de robustesse face aux conflits et aux indisponibilit\'es.
\end{itemize}

% =========================================
\section{Architecture du syst\`eme et interactions entre agents}
\subsection{Vue globale}
Le syst\`eme repose sur un \textbf{DispatcherAgent} jouant le r\^ole d'orchestrateur des interventions, entour\'e d'agents sp\'ecialis\'es (capteurs, unit\'es d'intervention, coordination hospitali\`ere, trafic, journalisation). Le d\'emarrage est s\'equenc\'e via \texttt{BootstrapAgent} pour assurer un enregistrement correct des services dans le DF (Directory Facilitator).

\begin{center}
\begin{tikzpicture}[
    node distance=0.9cm and 1.2cm,
    box/.style={draw, rounded corners, align=center, minimum width=3.1cm, minimum height=0.9cm},
    arr/.style={-{Latex[length=2mm]}, thick}
]
\node[box] (sensor) {Sensor Agents};
\node[box, right=of sensor] (dispatcher) {DispatcherAgent};
\node[box, above=of dispatcher] (traffic) {TrafficControllerAgent};
\node[box, below left=of dispatcher] (fire) {FireTruck Agents};
\node[box, below=of dispatcher] (amb) {Ambulance Agents};
\node[box, below right=of dispatcher] (bcu) {BCU Agents};
\node[box, right=of bcu] (police) {Police Agents};
\node[box, right=of amb] (med) {MedicalCoordinatorAgent};
\node[box, right=of med] (hosp) {Hospital Agents};
\node[box, above right=of dispatcher] (logger) {LoggerAgent};

\draw[arr] (sensor) -- node[above]{INFORM(HasEmergency)} (dispatcher);
\draw[arr] (dispatcher) -- node[right]{REQUEST route} (traffic);
\draw[arr] (dispatcher) -- (fire);
\draw[arr] (dispatcher) -- (amb);
\draw[arr] (dispatcher) -- (bcu);
\draw[arr] (dispatcher) -- (police);
\draw[arr] (amb) -- node[below]{REQUEST h\^opital} (med);
\draw[arr] (med) -- node[below]{affectation} (hosp);
\draw[arr] (dispatcher) -- (logger);
\end{tikzpicture}
\end{center}

\subsection{R\^oles et responsabilit\'es}
\begin{itemize}
    \item \textbf{SensorAgent} : d\'etection et publication d'incidents (message \texttt{INFORM} avec contenu ontologique \texttt{HasEmergency}).
    \item \textbf{DispatcherAgent} : qualification des incidents, lancement des appels d'offres, calcul d'utilit\'e, attribution des missions, relances et file d'attente.
    \item \textbf{FireTruckAgent / AmbulanceAgent / BiohazardContainmentUnitAgent / PoliceAgent} : ex\'ecution op\'erationnelle terrain avec cycle de mission.
    \item \textbf{TrafficControllerAgent} : gestion de demandes de d\'egagement de routes.
    \item \textbf{MedicalCoordinatorAgent + HospitalAgent} : s\'election d'h\^opital et gestion de capacit\'e en lits.
    \item \textbf{LoggerAgent} : tra\c{c}abilit\'e des \`ev\'enements et indicateurs de performance.
\end{itemize}

\subsection{Justification de l'architecture}
Cette architecture centralise la \textit{d\'ecision d'allocation} tout en distribuant l'\textit{ex\'ecution}. Ce compromis a \`et\'e retenu pour :
\begin{itemize}
    \item simplifier les politiques d'arbitrage multi-crit\`eres ;
    \item conserver des agents m\'etier autonomes et r\'eutilisables ;
    \item permettre l'extension par ajout de nouveaux services DF sans refonte globale.
\end{itemize}

% =========================================
\section{Protocoles utilis\'es : FIPA Contract Net et ontologie}
\subsection{Protocole FIPA Contract Net}
Le protocole principal suit les \`etapes suivantes :
\begin{enumerate}
    \item r\'eception d'un incident ;
    \item recherche DF des services compatibles (ex. \texttt{FIRE}, \texttt{MEDICAL}) ;
    \item envoi de \texttt{CFP} aux candidats ;
    \item r\'eception des \texttt{PROPOSE}/\texttt{REFUSE} ;
    \item \`a expiration du timeout : calcul des utilit\'es et choix du gagnant ;
    \item envoi \texttt{ACCEPT\_PROPOSAL} au gagnant, \texttt{REJECT\_PROPOSAL} aux autres.
\end{enumerate}

\textbf{Raison du choix :} Contract Net est adapt\'e aux environnements dynamiques et \`a ressources concurrentes, car il dissocie clairement publication de besoin et comp\'etition des r\'epondants.

\subsection{Protocole de messages compl\'ementaires}
En plus du Contract Net, le syst\`eme utilise des messages de cycle de vie :
\begin{itemize}
    \item \texttt{MISSION\_ARRIVED}, \texttt{MISSION\_COMPLETE}, \texttt{PERIMETER\_SECURED} ;
    \item \texttt{UNIT\_IDLE} pour signaler la disponibilit\'e ;
    \item \texttt{UNIT\_ABORT} / \texttt{ABORT} pour d\'eclencher une r\'eaffectation ;
    \item \texttt{HOSPITAL\_SATURATION} pour la gestion de congestion m\'edicale.
\end{itemize}

\subsection{Ontologie \texttt{Emergency-Ontology}}
L'ontologie formalise les objets m\'etier et les pr\'edicats de communication :
\begin{itemize}
    \item \textbf{Concepts} : \texttt{Location}, \texttt{Emergency}, \texttt{UnitStatus}, \texttt{Hospital}.
    \item \textbf{Pr\'edicats} : \texttt{HasEmergency}, \texttt{UnitAvailable}, \texttt{RouteCleared}, \texttt{HospitalAssignment}.
\end{itemize}

\textbf{Int\'er\^et de l'ontologie :}
\begin{itemize}
    \item coh\'erence s\'emantique des \`echanges ;
    \item meilleure robustesse qu'un protocole texte libre ;
    \item extensibilit\'e pour de nouveaux concepts (priorit\'e clinique, contraintes mat\'erielles, etc.).
\end{itemize}

% =========================================
\section{Fonction d'utilit\'e : math\'ematiques et logique du r\'epartiteur}
\subsection{D\'efinition}
Pour chaque unit\'e candidate $u$ et incident $e$, le r\'epartiteur calcule :
\[
U(u,e) = 0.5\,S_{type}(u,e)
+ 0.2\,\frac{1}{1+d(u,e)}
+ 0.15\,\frac{1}{1+w(u)}
+ 0.15\,\frac{sev(e)}{5}
\]

avec :
\begin{itemize}
    \item $S_{type}(u,e) \in [0,1]$ : compatibilit\'e de comp\'etence ;
    \item $d(u,e)$ : distance euclidienne entre position courante et lieu d'incident ;
    \item $w(u)$ : charge de l'unit\'e ;
    \item $sev(e) \in \{1,\dots,5\}$ : s\'ev\'erit\'e de l'incident.
\end{itemize}

\subsection{Justification des pond\'erations}
\begin{itemize}
    \item \textbf{0.50 sur le type} : priorit\'e \`a l'ad\'equation fonctionnelle (s\'ecurit\'e et pertinence op\'erationnelle).
    \item \textbf{0.20 sur la distance} : r\'eduction du temps d'arriv\'ee.
    \item \textbf{0.15 sur la charge} : \`equilibrage de l'utilisation des ressources.
    \item \textbf{0.15 sur la s\'ev\'erit\'e} : renforcement des situations critiques.
\end{itemize}

\subsection{R\`egles logiques compl\'ementaires}
La fonction num\'erique est compl\'et\'ee par des filtres discrets :
\begin{itemize}
    \item validation des \textit{primary responders} (ex. camion incendie pour \texttt{FIRE}) ;
    \item exclusion des unit\'es marqu\'ees occup\'ees ;
    \item alternative automatique si le meilleur candidat devient indisponible.
\end{itemize}

\textbf{Conclusion locale :} la d\'ecision finale est donc \`a la fois \textit{quantitative} (score) et \textit{symbolique} (contraintes m\'etier).

% =========================================
\section{Comportement des agents mod\'elis\'e en FSM}
\subsection{FSM du Dispatcher}
\begin{center}
\begin{tikzpicture}[node distance=1.5cm, >=Latex]
    \node[state,initial] (wait) {WAIT\_INCIDENT};
    \node[state,right=of wait] (cfp) {CALL\_FOR\_PROPOSALS};
    \node[state,right=of cfp] (collect) {COLLECT\_PROPOSE};
    \node[state,below=of collect] (eval) {EVALUATE};
    \node[state,left=of eval] (assign) {ASSIGN};
    \node[state,left=of assign] (retry) {RETRY/QUEUE};

    \path[->]
    (wait) edge node[above]{incident} (cfp)
    (cfp) edge node[above]{CFP envoy\'e} (collect)
    (collect) edge node[right]{timeout} (eval)
    (eval) edge node[below]{candidat valide} (assign)
    (eval) edge node[below]{\`echec} (retry)
    (assign) edge[bend left=18] node[above]{nouvel incident} (wait)
    (retry) edge[bend left=18] node[below]{relance} (cfp);
\end{tikzpicture}
\end{center}

\subsection{FSM des unit\'es d'intervention}
\textbf{AmbulanceAgent}
\begin{itemize}
    \item \texttt{IDLE} $\rightarrow$ \texttt{EN\_ROUTE} $\rightarrow$ \texttt{ON\_SITE} $\rightarrow$ \texttt{HOSPITAL\_TRANSPORT} $\rightarrow$ \texttt{RETURNING} $\rightarrow$ \texttt{IDLE}.
\end{itemize}

\textbf{FireTruckAgent}
\begin{itemize}
    \item \texttt{IDLE} $\rightarrow$ \texttt{EN\_ROUTE} $\rightarrow$ \texttt{ACTIVE} $\rightarrow$ \texttt{RETURNING}.
    \item branche de s\'ecurit\'e eau : \texttt{ACTIVE} $\rightarrow$ \texttt{REFILLING} si seuil critique.
\end{itemize}

\textbf{BiohazardContainmentUnitAgent}
\begin{itemize}
    \item \texttt{IDLE} $\rightarrow$ \texttt{EN\_ROUTE} $\rightarrow$ \texttt{ON\_SITE} $\rightarrow$ \texttt{DECONTAMINATING} $\rightarrow$ \texttt{RETURNING}.
\end{itemize}

\textbf{PoliceAgent}
\begin{itemize}
    \item \texttt{IDLE} (patrouille) $\rightarrow$ \texttt{EN\_ROUTE} $\rightarrow$ \texttt{ON\_SITE}.
    \item retour \`a \texttt{IDLE} sur \texttt{PERIMETER\_RELEASED}.
\end{itemize}

\subsection{Int\'er\^et de la mod\'elisation FSM}
La FSM permet :
\begin{itemize}
    \item une lisibilit\'e forte du comportement ;
    \item la pr\'evention des transitions incoh\'erentes ;
    \item une meilleure maintenabilit\'e et testabilit\'e des agents.
\end{itemize}

% =========================================
\section{Robustesse : gestion des deadlocks et des conflits}
\subsection{Conflits d'allocation et coh\'erence des ressources}
Le syst\`eme \`evite la double affectation par :
\begin{itemize}
    \item un ensemble global d'unit\'es occup\'ees (\texttt{busyUnits}) ;
    \item des validations au moment de l'acceptation ;
    \item des notifications \texttt{UNIT\_IDLE} pour lib\'eration explicite ;
    \item une recherche d'alternative imm\'ediate si le gagnant est indisponible.
\end{itemize}

\subsection{Pr\'evention des blocages (deadlocks fonctionnels)}
La conception JADE repose sur des comportements non bloquants :
\begin{itemize}
    \item \texttt{CyclicBehaviour} pour les r\'eceptions d'\`ev\'enements ;
    \item \texttt{WakerBehaviour} pour timeouts et transitions temporis\'ees ;
    \item \texttt{TickerBehaviour} pour surveillance p\'eriodique et reprise.
\end{itemize}

Cela \`evite les attentes actives et maintient la r\'eactivit\'e globale du syst\`eme.

\subsection{Tol\'erance aux pannes op\'erationnelles}
\begin{itemize}
    \item \textbf{ABORT mission} : ex. niveau d'eau insuffisant (camion incendie) ;
    \item \textbf{r\'eaffectation dynamique} : relance automatique d'un nouveau Contract Net ;
    \item \textbf{retries born\'es} : nombre maximal de relances ;
    \item \textbf{mise en file} : incidents non servis conserv\'es et retent\'es p\'eriodiquement ;
    \item \textbf{marquage FAILED} : fermeture propre des cas irr\'esolubles.
\end{itemize}

\subsection{Robustesse m\'edicale}
La couche hospitali\`ere traite la saturation :
\begin{itemize}
    \item affectation au centre le plus proche avec lit disponible ;
    \item signalement de saturation ;
    \item handoff patient avec r\'eponse explicite succ\`es/\'echec.
\end{itemize}

\subsection{Tra\c{c}abilit\'e et audit}
\texttt{LoggerAgent} enregistre les cycles d'incident (d\'etection, assignment, arriv\'ee, r\'esolution, \`echec), ce qui facilite :
\begin{itemize}
    \item le diagnostic de r\'egression ;
    \item l'\'evaluation de performance (temps de r\'eponse) ;
    \item la justification des choix de conception.
\end{itemize}

% =========================================
\section{Bilan critique}
\subsection{Points forts}
\begin{itemize}
    \item architecture modulaire et extensible ;
    \item d\'ecision explicable (score + contraintes m\'etier) ;
    \item bonne r\'esilience face aux indisponibilit\'es ;
    \item coh\'erence s\'emantique via ontologie partag\'ee.
\end{itemize}

\subsection{Limites observ\'ees}
\begin{itemize}
    \item pond\'erations d'utilit\'e statiques ;
    \item mod\`ele de trafic simplifi\'e ;
    \item centralisation partielle du calcul de d\'ecision.
\end{itemize}

\subsection{Perspectives}
\begin{itemize}
    \item adaptation dynamique des poids (approches data-driven) ;
    \item int\'egration d'indicateurs temporels forts (SLA) ;
    \item simulation de pannes r\'eseau et d\'efaillances simultan\'ees \`a grande \`echelle.
\end{itemize}

% =========================================
\section{Conclusion}
Le projet SRUU valide l'int\'er\^et d'une approche multi-agents pour la gestion d'urgences urbaines. L'association de \textbf{FIPA Contract Net}, d'une \textbf{ontologie m\'etier}, d'une \textbf{fonction d'utilit\'e multi-crit\`ere} et d'une \textbf{mod\'elisation FSM} fournit un syst\`eme coh\'erent, robuste et justifiable du point de vue ing\'enierie logicielle.

Les m\'ecanismes de reprise (timeouts, retries, queue, abort handling) r\'eduisent significativement les situations de blocage et am\'eliorent la continuit\'e de service. Cette base est suffisamment solide pour des extensions avanc\'ees orient\'ees optimisation et r\'esilience.

\end{document}
